\documentclass{beamer}
\usepackage{amsfonts,amsmath,oldgerm}
\usepackage{booktabs}
\usepackage{array}
\usepackage{graphicx}
\usepackage{tabularx}
\usepackage{xeCJK}
\usetheme{sintef}

\usefonttheme[onlymath]{serif}
\titlebackground*{assets/background}

\newcommand{\metric}[3]{\textsf{#1 / #2 / #3}}
\newcommand{\metriccell}[3]{\begin{tabular}{@{}c@{}}#1\\#2\\#3\end{tabular}}
\newcommand{\imgdir}{../images}
\newcommand{\smalltab}{\scriptsize}
\newcommand{\tight}{\vspace{-0.35em}}
\setlength{\tabcolsep}{3pt}
\renewcommand{\arraystretch}{1.08}

\title{实验课四：多模态大语言模型}
\subtitle{InternVL2-2B 监督微调、数据构造与显式 Grounding}
\author{group16 甜菜酱饭炼丹坊（田鑫芳 \ 蔡子晨 \ 姜证严 \ 范力文）}
\date{2024年5月19日}

\begin{document}
\maketitle

\begin{frame}{实验任务概览}
\small
\begin{columns}[T,onlytextwidth]
\begin{column}{0.53\textwidth}
\begin{block}{核心任务}
\begin{itemize}
\item 围绕 GQA VQA：输入 \texttt{image + question}，输出自然语言短答案。
\item 完整经历模型结构理解、SFT 训练、合成数据构造、评测与案例分析。
\end{itemize}
\end{block}
\end{column}
\begin{column}{0.43\textwidth}
\begin{block}{模型设置}
\begin{itemize}
\item 基座模型：\textbf{InternVL2-2B Base}
\item 教师模型：\textbf{Qwen3-VL-8B Instruct}
\end{itemize}
\end{block}
\end{column}
\end{columns}
\vspace{0.3em}
\begin{center}
\smalltab
\begin{tabular}{p{0.18\textwidth}p{0.39\textwidth}p{0.34\textwidth}}
\toprule
任务 & 可使用数据 & BBox 规则 \\
\midrule
Task 1 & GQA image/question/answer & 不允许使用 bbox \\
Task 2 & LVIS/COCO 非 QA 数据，可与 GQA 合并 & bbox 仅作构造中间信息，最终输入不能显式保留 \\
Task 3 & GQA、Task2 数据、GQA bbox & 允许显式用于训练或推理 \\
\bottomrule
\end{tabular}
\end{center}
\end{frame}

\section{Task 1：MLLM 微调策略探索}

\begin{frame}[shrink=4]{Baseline 构造与探索}
\small
\begin{itemize}
\item Task 1 只允许使用 GQA 的 \texttt{image + question + answer}，不使用 bbox 与 LVIS/COCO。
\item 早期实验发现：精确匹配评测不仅考察视觉理解，也高度依赖答案格式。
\item 推进至 Task 3 时发现，Prompt 对 InternVL2-2B 的影响显著，因此只能从 Task 1 起重新执行全部流程。
\end{itemize}
\vspace{0.2em}
\centering
\smalltab
\begin{tabular}{p{0.42\textwidth}ccc}
\toprule
设置 & Total & BBox & No-BBox \\
\midrule
InternVL2-2B, 原始 prompt & 0.247 & 0.2487 & 0.2222 \\
InternVL2-2B, 修正 prompt & \textbf{0.685} & 0.6756 & \textbf{0.8254} \\
Qwen3-VL-8B, 原始 prompt & \textbf{0.685} & \textbf{0.6777} & 0.7937 \\
Qwen3-VL-8B, 修正 prompt & 0.679 & 0.6692 & \textbf{0.8254} \\
\bottomrule
\end{tabular}
\vspace{0.4em}
\begin{block}{Prompt 修正}
\scriptsize
旧：\texttt{Question: \{question\}\textbackslash nAnswer with the shortest correct answer only.}\\
新：\texttt{Question: \{question\}\textbackslash nAnswer with the shortest correct answer only. The answer must be a simple short answer, usually one word, a short phrase, yes/no, or a number.}
\end{block}
\end{frame}

\begin{frame}{格式错误案例}
\small
\begin{columns}[T,onlytextwidth]
\begin{column}{0.45\textwidth}
\includegraphics[width=\textwidth,height=0.62\textheight,keepaspectratio]{\imgdir/task1_format_error_2386907.jpg}
\end{column}
\begin{column}{0.52\textwidth}
\begin{block}{样例}
\begin{itemize}
\item 图像：\texttt{images/gqa/2386907.jpg}
\item 问题：\textit{Is the grass brown or green?}
\item 标准答案：\texttt{green}
\item 模型输出：\texttt{The grass is green.}
\end{itemize}
\end{block}
\begin{block}{观察}
其实视觉理解是正确的，但精确匹配记为错误；因此后续统一使用修正 prompt。
\end{block}
\end{column}
\end{columns}
\end{frame}

\begin{frame}{模块冻结实验设计}
\small
\begin{itemize}
\item InternVL2-2B 包含视觉编码器、MLP Connector 和语言解码器。
\item 固定数据、prompt 与训练流程，仅改变可训练模块，比较准确率、效率和显存。
\end{itemize}
\vspace{0.2em}
\centering
\smalltab
\begin{tabular}{p{0.10\textwidth}p{0.31\textwidth}rrp{0.26\textwidth}}
\toprule
配置 & 可训练模块 & 参数量 & 占比 & 实验含义 \\
\midrule
A & Connector & 12.60M & 0.5710\% & 只学习跨模态对齐与格式适配 \\
B & Connector + Language & 1.90B & 86.2173\% & 冻结视觉编码器，标准 instruction tuning 风格 \\
C & Vision + Connector & 316.61M & 14.3537\% & 视觉表征适配，不更新语言模型 \\
D & Vision + Connector + Language & 2.21B & 100.0000\% & 全参数微调，作为性能上限候选 \\
\bottomrule
\end{tabular}
\vspace{0.5em}
\end{frame}

\begin{frame}{模块冻结实验：训练曲线}
\begin{columns}[T,onlytextwidth]
\begin{column}{0.76\textwidth}
\centering
\includegraphics[width=\linewidth,height=1.0\textheight,keepaspectratio]{\imgdir/task1_wandb.png}
\end{column}
\begin{column}{0.22\textwidth}
\scriptsize
\begin{block}{说明}
\begin{itemize}
\item 5090 上使用原始 prompt 训练。
\item 统一设置：\texttt{bs32}、\texttt{lr 1e-5}、\texttt{1024 samples}。
\item loss 很快接近 0，早期收益主要来自格式适配。
\item 控制相同 lr 在参数量差异巨大的情况下可能无法发挥最优性能。
\end{itemize}
\end{block}
\end{column}
\end{columns}
\end{frame}

\begin{frame}{实验结果：全参数最优，但成本收益非线性}
\small
\begin{itemize}
\item 模型：InternVL2-2B，no-bbox，修正 prompt
\item 统一设置：\texttt{batch size = 16}，\texttt{learning rate = 5e-6}，\texttt{512 samples}。
\end{itemize}
\centering
\smalltab
\begin{tabular}{p{0.20\textwidth}cccccc}
\toprule
配置 & Total & BBox & No-BBox & Trainable & Mem & Time \\
\midrule
Baseline & 0.685 & 0.6756 & 0.8254 & - & - & - \\
A Connector & 0.684 & 0.6756 & 0.8095 & 0.5710\% & 18.38G & 141.9s \\
B Conn.+Lang. & 0.701 & 0.6926 & 0.8254 & 86.2173\% & 19.39G & 176.4s \\
C Vision+Conn. & 0.688 & 0.6798 & 0.8095 & 14.3537\% & 20.41G & 229.5s \\
D Full & \textbf{0.707} & \textbf{0.6980} & \textbf{0.8413} & 100\% & 21.66G & 282.0s \\
\bottomrule
\end{tabular}
\vspace{0.3em}
\begin{block}{结论}
A 只训练 0.5710\% 参数小但作用有限；B 接近 D，说明语言侧 instruction tuning 对短答案 VQA 贡献稳定；C 未超过 B，在当前样本规模下单独更新视觉编码器并未展现出更优的效果。
\end{block}
\end{frame}

\begin{frame}[shrink=8]{Infra 与探索性实验}
\small
\begin{columns}[T,onlytextwidth]
\begin{column}{0.44\textwidth}
\begin{block}{Infra 修改}
\begin{itemize}
\item D\_full 开启 \texttt{gradient\_checkpointing}，避免 OOM。
\item 使用梯度累积灵活设置 batch size。
\item 图片预处理改为 CPU/GPU 并行，512 samples 从 297s 降至 282s。
\item 开启 Flash Attention 后降至 186s，共节省约 60\% 训练时间。
\end{itemize}
\end{block}
\end{column}
\begin{column}{0.54\textwidth}
\smalltab
\begin{tabular}{p{0.29\textwidth}rrp{0.35\textwidth}}
\toprule
设置 & Total & Time & 解释 \\
\midrule
 bs128 lr4e-5, 1024 & 0.588 & 588s & 官方参数，学习率偏大 \\
 bs64 lr2e-5, 1024 & 0.642 & 585s & 降低学习率后提升 \\
 bs32 lr2e-5, 1024 & 0.641 & 590s & 缩小 batch 收益有限 \\
 bs32 lr1e-5, 1024 & 0.691 & 590s & 进入较优学习率区间 \\
 bs32 lr5e-6, 1024 & 0.692 & 590s & 与 1e-5 接近，较稳定 \\
 bs32 lr1e-6, 1024 & 0.693 & 594s & 继续降学习率收益有限 \\
 bs16 lr5e-6, 1024 & 0.695 & 593s & 更小 batch 略有提升 \\
 bs16 lr5e-6, 512 & \textbf{0.699} & 297s & 样本减半不降反升 \\
 bs8 lr1e-6, 512 & 0.695 & 296s & 更小 batch 无明显优势 \\
 bs8 lr1e-6, 256 & 0.664 & 151s & 样本过少导致下降 \\
\bottomrule
\end{tabular}
\end{column}
\end{columns}
\vspace{0.3em}
\begin{block}{感悟与反思}
充分的 infra 与前期探索性实验很重要，搞好 infra 提升效率能更快获得反馈做更多消融实验，而 prompt 修正就是前期探索不足导致大批量返工的惨痛教训。
\end{block}
\end{frame}

\section{Task 2：从视觉标注数据构造 Instruction Tuning 数据}

\begin{frame}[shrink=5]{数据构造方法}
\small
\begin{itemize}
\item Task 2 使用 COCO caption 与 LVIS 物体类别、bbox 标注，将非 QA 数据转化为 instruction-response 样本。
\end{itemize}
\smalltab
\begin{tabular}{p{0.15\textwidth}p{0.28\textwidth}p{0.23\textwidth}p{0.24\textwidth}}
\toprule
方案 & 教师输入 & 输出类型 & 清洗重点 \\
\midrule
A text teacher & caption + objects + bbox 文本 & caption / 识别 / 计数 / 摘要 / 空间 & 去坐标泄漏、过滤非法 task type、去重 \\
B vision teacher & image + caption + objects + bbox 文本 & 同 A，但教师可看图纠偏 & 同 A，额外关注答案复杂度和视觉噪声 \\
\bottomrule
\end{tabular}
\vspace{0.4em}
\begin{block}{生成约束}
\begin{itemize}
\item 每张图最多生成 8 条样本，覆盖 \texttt{image\_captioning}、\texttt{object\_recognition}、\texttt{object\_counting}、\texttt{dense\_visual\_summary}、\texttt{spatial\_reasoning}。
\item 输出 JSON array，每条含 \texttt{task\_type/question/answer}；question 与 answer 使用英文。
\item 识别、计数和空间理解样本尽量做成短答案 VQA，适配精确匹配评测。
\end{itemize}
\end{block}
\end{frame}

\begin{frame}[fragile,shrink=4]{A/B 教师模型 Prompt 差异}
\small
\centering
\begin{minipage}{0.96\textwidth}
\begin{columns}[T,onlytextwidth,totalwidth=\linewidth]
\begin{column}{0.49\linewidth}
\vspace{0pt}
\centering
\begin{block}{方案 A：纯文字描述驱动}
\tiny
\begin{verbatim}
Objects in the image:
{fmt_objects(row, keep_bbox=True)}
Image caption: {row.get("caption", "")}
Please generate 8 diverse multimodal instruction-response samples.
Allowed task types:
image_captioning, object_recognition,
object_counting, dense_visual_summary, spatial_reasoning
\end{verbatim}
\end{block}
\end{column}
\begin{column}{0.49\linewidth}
\vspace{0pt}
\centering
\begin{block}{方案 B：图文联合驱动}
\tiny
\begin{verbatim}
Objects in the image:
{fmt_objects(row, keep_bbox=True)}
Image caption: {row.get("caption", "")}
Please look at the image and use the annotations as hints.
Generate 8 diverse multimodal instruction-response samples.
Allowed task types:
image_captioning, object_recognition,
object_counting, dense_visual_summary, spatial_reasoning
\end{verbatim}
\end{block}
\end{column}
\end{columns}
\vspace{-0.5em}
\begin{block}{共同规则}
\tiny
\begin{verbatim}
1. Output a JSON array. Each item must contain task_type, question, and answer.
2. task_type must be one of the allowed task types. Prefer diverse types.
3. question and answer must be in English.
4. For object_recognition/counting/spatial_reasoning, use VQA question +
   short answer: one word, phrase, yes/no, or number.
5. The final question must not contain bbox coordinates or bbox lists.
6. Output JSON only. Do not add explanations.
\end{verbatim}
\end{block}
\vspace{-0.5em}
\begin{block}{生成规模与分布}
\scriptsize
A：256 张图 $\times$ 8 $\rightarrow$ 2048 张保留 1789 张; B：同理保留 1522 张
\end{block}
\end{minipage}
\end{frame}

\begin{frame}[shrink=3]{数据质量案例一：118838}
\small
\begin{columns}[T,onlytextwidth]
\begin{column}{0.46\textwidth}
    \centering
    \includegraphics[width=\linewidth,height=1.0\textheight,keepaspectratio]{\imgdir/task2_case_118838_labeled.jpg}
\end{column}
\begin{column}{0.51\textwidth}
    \begin{block}{A text}
    \footnotesize
    \textit{How many gravestones are visible in the image?} $\rightarrow$ \texttt{3}\\[0.1em]
    仅凭标注生成计数。
    \end{block}
    \begin{block}{B vision}
    \footnotesize
    \textit{A bird sits on a gravestone, with two other gravestones and grassy surroundings visible.}\\[0.1em]
    看图后描述包含 \textit{grass / surroundings} 等额外视觉细节，信息更丰富。
    \end{block}
    \begin{block}{观察}
    \scriptsize
    A/B 的核心差异是教师模型是否拥有图像证据：B 的视觉细节更丰富，但也更可能偏向描述句而不是短答案 VQA。
    \end{block}
\end{column}
\end{columns}
\end{frame}

\begin{frame}[shrink=2]{数据质量案例二：388903}
\small
\begin{columns}[T,onlytextwidth]
\begin{column}{0.62\textwidth}
    \centering
    \includegraphics[width=\linewidth,height=1.0\textheight,keepaspectratio]{\imgdir/task2_case_388903_labeled.jpg}
\end{column}
\begin{column}{0.34\textwidth}
    \begin{block}{A text}
    \scriptsize
    \textit{How many suitcases are visible in the image?} $\rightarrow$ \texttt{1}\\[0.1em]
    结构化标注转写形式符合短答案要求，但图中其实不一定能看出来。
    \end{block}
    \vspace{-0.5em}
    \begin{block}{B vision}
    \scriptsize
    \textit{How many visible instances of the object 'apple' are there in the image?} $\rightarrow$ \texttt{many}\\[0.1em]
    看图后引入更真实的场景内容，但答案变成模糊量词，无法适配精确匹配评测。
    \end{block}
    \vspace{-0.5em}
    \begin{block}{结论}
    \scriptsize
    直接追求丰富 instruction 会引入 caption / summary 类数据，这些数据可能提升通用能力，但未必服务于 GQA 风格的短答案 VQA。
    \end{block}
\end{column}
\end{columns}
\end{frame}


\begin{frame}{Task 2 训练结果：没有稳定超过 GQA Baseline}
\small
\begin{itemize}
\item 模型：InternVL2-2B，D\_full 全参数微调。
\item 设置：\texttt{bs16 lr5e-6, 1024 samples}
\item 数据：单独和混合，混合时每条数据依次取自不同数据集并控制总样本数为 1024。
\end{itemize}
\smalltab
\centering 
\begin{tabular}{p{0.25\textwidth}ccccp{0.23\textwidth}}
\toprule
训练数据 & Total & BBox & No-BBox & $\Delta$ & 解释 \\
\midrule
Baseline & \textbf{0.707} & \textbf{0.6980} & \textbf{0.8413} & 0.000 & GQA 原生监督最强 \\
A only & 0.691 & 0.6830 & 0.8095 & -0.016 & 合成数据不能替代 GQA \\
B only & 0.687 & 0.6788 & 0.8095 & -0.020 & 比 A 还差 \\
A+B & 0.691 & 0.6830 & 0.8095 & -0.016 & 与 A 基本相同 \\
Task1+A & 0.697 & 0.6884 & 0.8254 & -0.010 & A 混合相比纯A略提升 \\
Task1+B & 0.699 & 0.6905 & 0.8254 & -0.008 & B 混合提升更大 \\
Task1+A+B & 0.695 & 0.6862 & 0.8254 & -0.012 & 多源噪声累积 \\
\bottomrule
\end{tabular}
\end{frame}

\begin{frame}{Task 2 结论：主要瓶颈是分布偏差}
\small
\begin{columns}[T,onlytextwidth]
\begin{column}{0.48\textwidth}
\begin{block}{正向结果}
\begin{itemize}
\item 非 QA 标注可批量生成 instruction-response。
\item 数据覆盖 caption、识别、计数、摘要和空间关系。
\item B 方案能引入更多真实视觉细节。
\end{itemize}
\end{block}
\end{column}
\begin{column}{0.48\textwidth}
\begin{block}{问题}
\begin{itemize}
\item 合成样本与 GQA 短答案 VQA 的问题风格不同。
\item 答案粒度、视觉 grounding 可靠性与验证集不一致。
\item 追求丰富 instruction 会引入 caption/summary 类分布偏差。
\end{itemize}
\end{block}
\end{column}
\end{columns}
\vspace{0.4em}
\begin{block}{后续尝试}
后续又生成了更偏 VQA 短回答风格的数据，但仍没有有效提升；因此 Task 3 转向显式使用 bbox grounding，而不是继续增加弱监督合成 QA。
\end{block}
\end{frame}

\section{Task 3：Bounding Box 信息的利用}

\begin{frame}[shrink=4]{从隐式视觉理解到显式 Grounding}
\small
\begin{itemize}
\item Task 1/2 的共同瓶颈是定位敏感题：模型需要先找到正确对象，再回答相关问题。
\item Task 3 允许显式使用 GQA bbox，问题变为如何更好地利用 grounding 信息。
\item 验证集中 \textbf{937/1000} 个样本含 bbox，预期收益很大。
\end{itemize}
\begin{columns}[T,onlytextwidth]
\begin{column}{0.42\textwidth}
\centering
\begin{tabular}{rrr}
\toprule
BBox 数量 & 样本数 & 占比 \\
\midrule
0 & 63 & 6.3\% \\
1 & 394 & 39.4\% \\
2 & 451 & 45.1\% \\
3 & 91 & 9.1\% \\
4 & 1 & 0.1\% \\
\bottomrule
\end{tabular}
\end{column}
\begin{column}{0.55\textwidth}
\vspace{-1em}
\begin{block}{探索方案}
\begin{itemize}
\item 文本化 bbox：类别 + 坐标拼接到问题前。
\item 图像可视化：同色框、彩色框、彩色框 + 标签。
\item bbox-aware CoT：参考 VoCoT\footnote{\href{https://arxiv.org/abs/2405.16919}{VoCoT: Unleashing Visually Grounded Multi-Step Reasoning in Large Multi-Modal Models}} 探索推理。
\item hard case：基于训练集错误样本再训练。
\item 基于 BBox 裁切区域图像方案在 VoCoT 中证明无效，没有做更多尝试。
\end{itemize}
\end{block}
\end{column}
\end{columns}
\end{frame}

\begin{frame}[fragile]{Prompt 侧方案}
\small
\begin{columns}[T,onlytextwidth]
\begin{column}{0.49\textwidth}
\begin{block}{no\_bbox（无信息）}
\tiny
\begin{verbatim}
Question:
Is the mattress to the right or to the left
of the table the lamp is on?
Answer with the shortest correct answer only.
\end{verbatim}
\end{block}
\begin{block}{color\_prompt（框颜色+标签）}
\tiny
\begin{verbatim}
Objects in the image:
1. red box: lamp
2. blue box: table
3. green box: mattress

Question: Is the mattress to the right or to
the left of the table the lamp is on?
Answer with the shortest correct answer only.
\end{verbatim}
\end{block}
\end{column}
\begin{column}{0.49\textwidth}
\begin{block}{bbox\_prompt（标签+定位）}
\tiny
\begin{verbatim}
Objects in the image:
1. lamp: normalized_bbox=[0.172, ..., 0.661064]
2. table: normalized_bbox=[0.174, ..., 0.935574]
3. mattress: normalized_bbox=[0.348, ..., 0.966387]

Question: Is the mattress to the right or to
the left of the table the lamp is on?
Answer with the shortest correct answer only.
\end{verbatim}
\end{block}
\begin{block}{color\_bbox\_prompt（三种都有）}
\tiny
\begin{verbatim}
Objects in the image:
1. red box: lamp, normalized_bbox=[...]
2. blue box: table, normalized_bbox=[...]
3. green box: mattress, normalized_bbox=[...]

Question: Is the mattress to the right or to
the left of the table the lamp is on?
Answer with the shortest correct answer only.
\end{verbatim}
\end{block}
\end{column}
\end{columns}
\end{frame}

\begin{frame}{图像侧方案：bbox 可视化}
\small
\begin{columns}[T,onlytextwidth]
\begin{column}{0.32\textwidth}
\includegraphics[width=\textwidth,height=0.48\textheight,keepaspectratio]{\imgdir/task3_0953938_same_box.jpg}
\centering\scriptsize \texttt{same\_box}\\
同色框提示区域，但难区分多目标
\end{column}
\begin{column}{0.32\textwidth}
\includegraphics[width=\textwidth,height=0.48\textheight,keepaspectratio]{\imgdir/task3_0953938_color_box.jpg}
\centering\scriptsize \texttt{color\_box}\\
彩色框支持颜色引用和多目标区分
\end{column}
\begin{column}{0.32\textwidth}
\includegraphics[width=\textwidth,height=0.48\textheight,keepaspectratio]{\imgdir/task3_0953938_color_box_label.jpg}
\centering\scriptsize \texttt{color\_box\_label}\\
彩色框 + 编号/类别标签
\end{column}
\end{columns}
\vspace{0.4em}
\begin{block}{设计取舍}
可视化把位置信息变成视觉提示，但可能遮挡原图或引入噪声；标签往往比框本身更关键。
\end{block}
\end{frame}

\begin{frame}[shrink=8]{Prompt × 图像二维消融}
\small
\begin{itemize}
\item 不训练InternVL2-2B，直接比较不同推理输入形式，判断 bbox 信息形式本身是否有效。
\item 每个单元为 \texttt{Total / BBox / No-BBox}；最稳定的是 \texttt{bbox\_prompt + original image}。
\item 标签/文本上下文是核心：\texttt{color\_box\_label} 在 \texttt{no\_bbox} 下有效；已有 bbox prompt 时，额外画框通常不增益，甚至造成视觉噪声。
\item color\_prompt + original image 不带任何定位信息，只有类别文本，甚至还有颜色噪声，但表现意外地也很好，是因为给出了类别文本，很多答案文本就是某个类别，避免了同义词错误。
\item no\_bbox + color box label 相比之下提升很大，说明图上文本也有帮助。
\end{itemize}

% 关键：表格全局居中
\centering
\scriptsize
\setlength{\tabcolsep}{2.5pt}
\resizebox{\textwidth}{!}{%
\begin{tabular}{lcccc}
\toprule
Prompt \textbackslash Image & original & \shortstack{same box} & \shortstack{color box} & \shortstack{color box label} \\
\midrule
no\_bbox & \metric{0.685}{0.6756}{\textbf{0.8254}} & \metric{0.668}{0.6574}{\textbf{0.8254}} & \metric{0.670}{0.6596}{\textbf{0.8254}} & \metric{0.796}{0.7940}{\textbf{0.8254}} \\
bbox\_prompt & \metric{\textbf{0.799}}{\textbf{0.7972}}{\textbf{0.8254}} & \metric{0.781}{0.7780}{\textbf{0.8254}} & \metric{0.781}{0.7780}{\textbf{0.8254}} & \metric{0.790}{0.7876}{\textbf{0.8254}} \\
color\_prompt & \metric{0.798}{0.7962}{\textbf{0.8254}} & \metric{0.782}{0.7791}{\textbf{0.8254}} & \metric{0.772}{0.7684}{\textbf{0.8254}} & \metric{0.779}{0.7759}{\textbf{0.8254}} \\
color\_bbox & \metric{0.798}{0.7962}{\textbf{0.8254}} & \metric{0.758}{0.7535}{\textbf{0.8254}} & \metric{0.756}{0.7513}{\textbf{0.8254}} & \metric{0.790}{0.7876}{\textbf{0.8254}} \\
\bottomrule
\end{tabular}
}
\vspace{0.4em}
\end{frame}


\begin{frame}[shrink=8]{训练方式探索}
\small
\begin{itemize}
\item Prompt × 图像二维消融说明单纯bbox prompt效果最优，但训练后有可能被反超吗？
\item 超参数统一使用 512 samples + bs16 + lr5e-6，实验名为 prompt + 图片 方案
\end{itemize}
\vspace{0.15em}
\scriptsize
\centering 
\begin{tabular}{p{0.36\textwidth}cccp{0.29\textwidth}}
\toprule
实验 & Total & BBox & No-BBox & 解释 \\
\midrule
no\_bbox 训练, bbox\_prompt 推理 & 0.832 & 0.8324 & 0.8254 & 仅推理加 bbox 已明显超过 Task1 \\
bbox\_prompt + orginal & 0.867 & 0.8687 & \textbf{0.8413} & 仅 prompt 加 prompt 效果不错 \\
no\_bbox + same\_box & 0.681 & 0.6702 & \textbf{0.8413} & 纯视觉同色框不太行 \\
no\_bbox + color\_box & 0.686 & 0.6756 & \textbf{0.8413} & 彩色框也不太行 \\
no\_bbox + color\_box\_label & 0.822 & 0.8218 & 0.8254 & 图上标签带来明显改善 \\
bbox\_prompt + same\_box & 0.860 & 0.8613 & \textbf{0.8413} & 文本 bbox 是核心 \\
color\_prompt + color\_box & 0.852 & 0.8527 & \textbf{0.8413} & 颜色 prompt 有效但低于 bbox prompt \\
color\_box\_prompt + color\_box & 0.860 & 0.8613 & \textbf{0.8413} & prompt 中的坐标比画框更关键 \\
bbox\_prompt + 4096 samples, bs64 & \textbf{0.874} & \textbf{0.8762} & \textbf{0.8413} & 扩大规模有收益 \\
\bottomrule
\end{tabular}
\vspace{0.4em}
\begin{block}{结论}
训练与推理阶段的 bbox 表达一致性很重要；单独画框不稳定，还是得简单直接的 bbox prompt。
\end{block}
\end{frame}


\begin{frame}[shrink=4]{教师模型参考}
\small
\begin{center}
\smalltab
\begin{tabular}{p{0.43\textwidth}cccp{0.22\textwidth}}
\toprule
模型与 prompt + 图像方案 & Total & BBox & No-BBox & 观察 \\
\midrule
InternVL2-2B, no\_bbox + original & 0.685 & 0.6756 & 0.8254 & 作为基础对照 \\
InternVL2-2B, bbox\_prompt + original & 0.799 & 0.7972 & 0.8254 & 显式 bbox 已显著增益 \\
InternVL2-2B, color\_bbox\_prompt + color\_box\_label & 0.790 & 0.7876 & 0.8254 & 画框未超过纯 bbox prompt \\
Qwen3-VL-8B, no\_bbox + original & 0.679 & 0.6692 & 0.8254 & 强模型也受定位限制 \\
Qwen3-VL-8B, bbox\_prompt + original & 0.842 & 0.8431 & 0.8254 & 大幅提升 \\
Qwen3-VL-8B, color\_bbox\_prompt + color\_box\_label & \textbf{0.844} & \textbf{0.8453} & 0.8254 & 图上信息仍有增益 \\
\bottomrule
\end{tabular}
\end{center}
\begin{block}{启发}
两种模型都能从 bbox 信息中获益；但 Intern 在信息过多时表现反而下降，而 Qwen 则能继续提升。然而 Qwen 的最优结果也没有超过探索试验中 Intern 微调后的最优结果，枉为人师，因此后续均只使用 Intern。
\end{block}
\end{frame}


\begin{frame}[shrink=6]{bbox\_prompt训练：模块冻结结果}
\small
\begin{itemize}
\item 在 bbox prompt 推理有效的基础上，训练目标变为让模型稳定适配输入格式。
\item 本页统一使用 bbox\_prompt 输入格式，训练样本数固定为 512。
\item 学习率使用各配置参数量较优值：A=1e-4，B=5e-6，C=1e-5，D=5e-6。
\item \textbf{D\_full 最高，但 B\_connector\_language 非常接近；语言侧对 bbox 文本格式的适配更重要，但视觉侧对 no-bbox 有效。}
\end{itemize}
\centering  % 表格居中    % 表格缩小一点，避免溢出
\begin{tabular}{lccccccp{0.18\textwidth}}
\toprule
配置 & Total & BBox & No-BBox & Trainable & Mem & Time & 观察 \\
\midrule
Baseline & 0.799 & 0.7972 & 0.8254 & - & - & - & 未训练 \\
A Connector & 0.855 & 0.8581 & 0.8095 & 0.5710\% & 19.17G & 274s & 轻量对齐也能利用 bbox \\
B Conn.+Lang. & 0.869 & 0.8719 & 0.8254 & 86.2173\% & 19.36G & 376s & 语言侧适配关键 \\
C Vision+Conn. & 0.853 & 0.8538 & 0.8413 & 14.3537\% & 21.20G & 499s & 视觉侧收益较小 \\
D Full & \textbf{0.872} & \textbf{0.8751} & 0.8254 & 100\% & 21.63G & 603s & 后续强基线 \\
\bottomrule
\end{tabular}
\vspace{0.3em}
\end{frame}


\begin{frame}[fragile]{CoT 与 Hard Case 尝试}
\small
\begin{columns}[T,onlytextwidth]
\begin{column}{0.50\textwidth}
\begin{block}{COT 结果}
\begin{itemize}
\item 参考 VoCoT，让模型先基于 bbox 生成区域描述和推理路径，再输出短答案。
\item Qwen 在 100 个样本上：直接 \texttt{bbox\_prompt} 为 \metric{0.840}{0.8404}{0.8333}。
\item CoT 为 \metric{0.810}{0.8085}{0.8333}，无法合成更优质数据。
\end{itemize}
\scriptsize
格式：
\vspace{-1em}
\begin{verbatim}
Thought: <brief reasoning process>
Answer: <short answer>
\end{verbatim}
\end{block}
\end{column}
\begin{column}{0.47\textwidth}
\begin{block}{Hard Case}
\begin{itemize}
\item 用 \texttt{bbox\_prompt} 跑 8000 条训练数据。
\item 正确 6436/8000 = 0.8045，筛出 1564 条错误样本。
\item 只训 hard cases 后验证 \metric{0.834}{0.8346}{0.8254}，不如同等规模训普通样本。
\item 后续改为 hard cases 与常规 Task1/Task3 样本混合。
\end{itemize}
\end{block}
\end{column}
\end{columns}
\end{frame}

\begin{frame}{最终方案：bbox prompt + hard case 混合训练}
\small
\begin{itemize}
\item 第一阶段：使用 \texttt{bbox\_prompt}，D\_full，在 task3 8000 条数据上训练，lr 使用 0.05 warmup + cosine 衰减。
\item 第二阶段：原始模型在 task1/3 训练集上的 hard cases，针对性提高难题能力。
\item 第三阶段：再在打乱的 task3 8000 条数据上训练恢复一下分布，取验证集最优。
\item 最终提交 Leaderboard 得分 \textbf{85.3}，排名第一。
\end{itemize}
\centering 
\smalltab
\begin{tabular}{p{0.42\textwidth}cccp{0.21\textwidth}}
\toprule
训练设置 & Total & BBox & No-BBox & 观察 \\
\midrule
baseline：bbox\_prompt & 0.799 & 0.7972 & 0.8254 & 未训练模型 \\
task3 8000 samples, bs16 lr1e-5 & 0.890 & 0.8922 & 0.8571 & 长训练大幅提升 \\
task1/3 hard 3128 samples, bs16 lr5e-6 & 0.894 & 0.8954 & 0.8730 & hard case 混合有效 \\
task3 shuffle 8000 samples, bs64 lr5e-6 / best & \textbf{0.906} & \textbf{0.9082} & \textbf{0.8730} & step 100 最优 checkpoint \\
\bottomrule
\end{tabular}
\vspace{0.35em}
\end{frame}

\begin{frame}{最终训练曲线与分析}
\small
\begin{columns}[T,onlytextwidth]
% 左列调窄，图片右移+放大
\begin{column}{0.6\textwidth}
% 调大尺寸，消除留白空隙
\includegraphics[width=1.0\textwidth,height=1.0\textheight,keepaspectratio]{\imgdir/task3_wandb.png}
\end{column}
% 右列同步加宽，贴合排版
\begin{column}{0.38\textwidth}
\begin{block}{关键结论}
\begin{itemize}
\item Task 3 的主要提升来自显式 grounding，而不是更复杂的 instruction 数据。
\item bbox prompt 把定位问题转化为模型可读的结构化上下文。
\item 标签/对象文本比单纯可视化画框更有效。
\item CoT 对当前简单短答案精确匹配任务收益不大。
\end{itemize}
\end{block}
\end{column}
\end{columns}
\end{frame}

\section{总结}

\begin{frame}[shrink=4]{总体结论}
\footnotesize
\begin{columns}[T,onlytextwidth]
\begin{column}{0.49\textwidth}
\begin{block}{Task 1}
\begin{itemize}
\item Prompt 格式约束会显著影响 VQA 精确匹配结果。
\item 当前规模下语言侧适配比视觉编码器微调更有效。
\item Connector 参数少但性价比高。
\end{itemize}
\end{block}
\begin{block}{Task 2}
\begin{itemize}
\item 从 LVIS/COCO 标注构造 instruction 数据是可行的。
\item B 方案细节更丰富但噪声更大；A 方案更保守，更贴近短答案格式。
\item 仅用合成数据或与 GQA 混合训练都未带来稳定正增益。
\end{itemize}
\end{block}
\end{column}
\begin{column}{0.49\textwidth}
\begin{block}{Task 3}
\begin{itemize}
\item 显式 bbox grounding 是最直接的增益来源，且文本大于图像。
\item CoT 和只训 hard cases 都不稳定。
\item 最终简单的训练方案就能达到 0.906 验证准确率。
\end{itemize}
\end{block}
\begin{block}{反思与感悟}
\begin{itemize}
\item 先把 prompt、infra、评测、可观测性等等搞好，不要盲目乱训。
\item 合成数据与目标分布的一致性，比单纯增加数据量更关键。
\item 谨记奥卡姆剃刀，大道至简，但也需要充分系统的消融实验支撑。
\end{itemize}
\end{block}
\end{column}
\end{columns}
\end{frame}

\backmatter
\end{document}
